\clearpage
\fancyhead[L]{\small GAUSS AN BODE.}
\hypertarget{native-444-part-lower}{}
\begin{center}
{\large GAUSS AN BODE.}

\noindent\rule{\linewidth}{0.4pt}

Astronomisches Jahrbuch für 1825. S. 103..105. Berlin 1821.

\noindent\rule{\linewidth}{0.4pt}
\end{center}

\begin{flushright}
Göttingen, den 26. Dezember 1821.
\end{flushright}

Sie werden mir mein langes Schweigen nachsehen, wenn ich Ihnen sage, dass ich den größten Teil des Jahres von hier fort war und selbst noch im Spätherbst nach Altona gereist bin, um den Zenitsektor von \textsc{Ramsden} in Empfang zu nehmen. Von dort bin ich erst vor Kurzem zurückgekehrt.

Bei meiner Triangulation, für die ich bisher an fünf Eckpunkten die Winkel gemessen habe, habe ich versucht, die Dreiecke so groß wie möglich zu machen. Sie werden in Nr. 126 der hiesigen Göttingischen gelehrten Anzeigen [Bd.\,IV der Werke.] die Mitteilung über das neue Hilfsmittel gelesen haben, das ich dafür einsetze, den \emph{Heliotrop}, und über die ersten Versuche im großen Maßstab, die damit gemacht wurden. Seitdem habe ich es ständig benutzt, als Zielpunkt beim Winkelmessen und mit ebenso gutem Erfolg für telegrafische Signale. Das Sonnenlicht, das ein Spiegel von 2 Zoll Breite und $1\tfrac{1}{4}$ Zoll Höhe zurückwarf, war aus Entfernungen von 5, 6, $7\tfrac{1}{2}$, ja einmal von $9\tfrac{1}{3}$ geografischen Meilen mit \emph{bloßem Auge} zu sehen. Seine gewaltige Wirkung versetzt gewöhnlich diejenigen in Staunen, die sie zum ersten Mal erleben und durch keine theoretische Berechnung darauf vorbereitet sind. Bei einigermaßen günstigen Luftverhältnissen sind der Länge der Dreiecksseiten jetzt nur noch die Grenzen gesetzt, die sich aus der Krümmung der Erde ergeben. Das gilt besonders, wenn man die Spiegel noch etwas größer macht, wie ich es bei zwei neu angefertigten Heliotropen ganz verschiedener Bauart getan habe.

Die genannten Arbeiten haben mich bisher daran gehindert, mich um die Berechnung einer Ephemeride für die nächste Erscheinung der Pallas zu kümmern und die früheren Beobachtungen, die ich hier gemacht habe, für Ihr Jahrbuch zu ordnen. Fürs Erste kann ich Ihnen hier nur meine Beobachtungen des diesjährigen Kometen beifügen.

\Needspace{15\baselineskip}
\begin{center}
\small
\setlength{\tabcolsep}{4pt}
\begin{tabular}{l|rrr|rrr|rrr}
1821 & \multicolumn{3}{c|}{\begin{tabular}{c}Mittlere Zeit\\ in Göttingen\end{tabular}} & \multicolumn{3}{c|}{Rektaszension} & \multicolumn{3}{c}{Nördliche Deklination} \\
\hline
Jan. 30 & 7\textsuperscript{u} & 34\textsuperscript{m} & 42\textsuperscript{s} & $359^\circ$ & $27'$ & $7''$ & $16^\circ$ & $4'$ & $36''$ \\
Febr. 3 & 7 & 3 & 56 & 359 & 3 & 54 & 15 & 46 & 3 \\
 7 & 6 & 42 & 38 & 358 & 45 & 5 & 15 & 29 & 49 \\
 9 & 6 & 42 & 30 & 358 & 36 & 24 & 15 & 21 & 22 \\
 10 & 6 & 52 & 27 & 358 & 32 & 19 & 16 & 17 & 50 \\
 11 & 7 & 12 & 0 & 358 & 28 & 27 & 15 & 14 & 26 \\
März 1 & 7 & 18 & 7 & 357 & 13 & 34 & 14 & 8 & 37 \\
 5 & 7 & 5 & 28 & 356 & 54 & 11 & 13 & 43 & 5 \\
\end{tabular}
\end{center}

\Needspace{18\baselineskip}
\SourceLeaf{445}
Die parabolischen Elemente, die Herr von \textsc{Staudt} daraus berechnet hat, sind folgende.

\begin{center}
\small
\begin{tabular}{p{.43\linewidth}p{.49\linewidth}}
Zeit des Perihels \mbox{.\,.\,.\,.\,.\,.\,.} & 1821 März $21.61890$ in Göttingen \\
—\quad Arg. der Breite \mbox{.\,.\,.\,.\,.\,.} & $169^\circ\;10'\;9''3$ \\
—\quad Logarithmus des Abstands \mbox{.\,.\,.} & $8.9645990$ \\
{\rightskip=0pt plus 3em Länge des aufsteigenden Knotens \mbox{.\,.\,.\,.\,.}\par} & $48^\circ\;44'\;14''7$ vom mittl. Äquinoktium des 1. Febr. \\
Neigung der Bahn \mbox{.\,.\,.\,.\,.\,.\,.\,.\,.\,.} & $106\quad40\quad16.4$ \\
\end{tabular}
\end{center}

Herr von \textsc{Staudt} hat diese Elemente sowohl mit meinen Beobachtungen als auch mit allen anderen verglichen, die damals bekannt geworden waren. Die Ergebnisse dieses Vergleichs lege ich für Ihr Jahrbuch in seiner eigenen Handschrift bei, weil die Zeit zum Abschreiben zu knapp ist.

\Needspace{34\baselineskip}
\begin{center}
\begin{tabular}{l|r|r}
& \multicolumn{2}{c}{Differenzen in} \\
& AR. & Dekl. \\
\multicolumn{3}{c}{Paris} \\
Jan. 21 & $+0''2$ & $+0''1$ \\
\multicolumn{3}{c}{Göttingen} \\
Jan. 30 & $-18''8$ & $+34''9$ \\
Febr. 3 & $+3.9$ & $-6.0$ \\
7 & $-13.1$ & $-36.0\!:$ \\
9 & $-0.1$ & $+14.7$ \\
10 & $+2.6$ & $+7.3$ \\
11 & $-1.2$ & $-3.5$ \\
März 1 & $-0.3$ & $-0.3$ \\
5 & $-5.3$ & $-20.6$ \\
\multicolumn{3}{c}{Bremen} \\
Jan. 30 & $-12.6$ & $+12.4$ \\
30 & $+9.1$ & $+34.2$ \\
Febr. 2 & $+24.7$ & $+9.7$ \\
5 & $-7.3$ & $-42.3\!::$ \\
7 & $+8.8$ & $+16.0\!:$ \\
8 & $+4.9$ & $+22.4$ \\
9 & $+5.0$ & $+14.1$ \\
10 & $-5.8$ & $+20.2$ \\
11 & $+4.4$ & $+3.2$ \\
12 & $-9.5$ & $+0.4$ \\
13 & $+0.6$ & $-24.6$ \\
14 & $-2.8$ & $-19.6$ \\
19 & $+10.2$ & $-3.4$ \\
März 1 & $+8.5$ & $-9.4$ \\
5 & $-0.6$ & $-7.6$ \\
6 & $+8.6$ & $-0.2$ \\
\end{tabular}
\end{center}

\Needspace{38\baselineskip}
\begin{center}
\begin{tabular}{l|r|r}
& \multicolumn{2}{c}{Differenzen in} \\
& AR. & Dekl. \\
\multicolumn{3}{c}{Mannheim} \\
Febr. 6 & $-0''2$ & $+28''9\!:$ \\
7 & $+2.2$ & $+1.3$ \\
8 & $+9.6$ & $-1.1$ \\
9 & $-6.2$ & $+5.7$ \\
10 & $+5.9$ & $+7.3$ \\
11 & $+12.0$ & $+17.4$ \\
12 & $+8.6$ & $+6.3$ \\
13 & $+8.6$ & $+4.2$ \\
14 & $+9.2$ & $+12.7$ \\
15 & $+9.4$ & $-3.7$ \\
27 & $+3.7$ & $-18.2$ \\
\multicolumn{3}{c}{Königsberg} \\
Febr. 9 & $-1.2$ & $+0.5$ \\
10 & $-5.4$ & $+2.9$ \\
11 & $+1.2$ & $-7.6$ \\
12 & $+1.9$ & $-1.1$ \\
14 & $+14.3$ & $-6.3$ \\
15 & $+11.7$ & $-20.8$ \\
19 & $-9.5$ & $+20.7$ \\
25 & $+35.3$ & $+35.9$ \\
27 & $+5.1$ & $-10.6$ \\
März 4 & $+15.6$ & $-12.4$ \\
5 & $-11.4$ & $+4.3$ \\
6 & $+7.8$ & $-17.2$ \\
\multicolumn{3}{c}{Hamburg} \\
Febr. 7 & $-3.2$ & $-29.2$ \\
8 & $+4.5$ & $+28.9$ \\
9 & $-6.5$ & $+9.9$ \\
10 & $+4.5$ & \mbox{.\,.\,.\,.} \\
12 & $-7.0$ & \mbox{.\,.\,.\,.} \\
\multicolumn{3}{c}{Dorpat} \\
Febr. 14 & $-3.6$ & $+83.0\!::$ \\
\end{tabular}
\end{center}

\Needspace{39\baselineskip}
\begin{center}
\begin{tabular}{l|r|r}
& \multicolumn{2}{c}{Differenzen in} \\
& AR. & Dekl. \\
\multicolumn{3}{c}{Seeberg} \\
Febr. 3 & $+30''2$ & $+103''4\!::$ \\
5 & $-8.6$ & $-41.5$ \\
7 & $+14.3$ & $-95.1\!::$ \\
8 & $+25.1$ & $+40.3\!::$ \\
9 & $+9.8$ & $-20.6$ \\
10 & $-4.9$ & $-19.0$ \\
11 & $+0.1$ & $-15.1$ \\
12 & $+10.2$ & $-20.0$ \\
14 & $+5.3$ & $-39.7$ \\
\multicolumn{3}{c}{Wien} \\
Febr. 20 & $-78.8$ & $-28.1$ \\
\multicolumn{3}{c}{Mailand} \\
Jan. 31 & $-37.7$ & $+4.5$ \\
Febr. 1 & $-22.9$ & $-4.4$ \\
2 & $-0.2$ & $-3.1$ \\
3 & $-29.4$ & $-0.2$ \\
7 & $-31.5$ & $-10.5$ \\
8 & $-19.8$ & $-2.3$ \\
\multicolumn{3}{c}{Padua} \\
Febr. 16 & $-50.5$ & $-56.1$ \\
16 & $-30.8$ & $-60.2$ \\
17 & $-69.4$ & $-15.5$ \\
17 & $-44.0$ & $-45.9$ \\
18 & $-61.5$ & $-1.3$ \\
18 & $-34.6$ & $-30.1$ \\
\multicolumn{3}{c}{Florenz} \\
Febr. 16 & $-14.8$ & $-24.9$ \\
17 & $-7.0$ & $-12.9$ \\
21 & $+5.9$ & $-47.6$ \\
22 & $-19.1$ & $-34.8$ \\
23 & $-8.6$ & $-47.5$ \\
25 & $-26.3$ & $-265.0$ \\
27 & $+16.0$ & $-603.2$ \\
\end{tabular}
\end{center}
